% !TEX program = pdflatex
% !BIB program = biber
\documentclass[oneside,final,11pt,letterpaper]{article}
% Shared formatting for course reports in this workspace.
% Packages, bibliography and the one-page title. Size and spacing for a
% specific course go in that course's preamble (MSD: 12 pt, 5-page body).
% Each assignment: \documentclass[...]{article} then % Shared formatting for course reports in this workspace.
% Packages, bibliography and the one-page title. Size and spacing for a
% specific course go in that course's preamble (MSD: 12 pt, 5-page body).
% Each assignment: \documentclass[...]{article} then % Shared formatting for course reports in this workspace.
% Packages, bibliography and the one-page title. Size and spacing for a
% specific course go in that course's preamble (MSD: 12 pt, 5-page body).
% Each assignment: \documentclass[...]{article} then \input{../../common/preamble}

\usepackage{tempora} % Times New Roman alike font

\usepackage{geometry}
\geometry{left=2.2cm, right=2.2cm, top=1.8cm, bottom=1.8cm, headheight=15pt, headsep=6mm, footskip=10mm}
\usepackage{setspace}
\setstretch{1.15}
\usepackage{indentfirst}
\parindent=1em
\setlength{\parskip}{0pt}

\usepackage[T1]{fontenc}
\usepackage[utf8]{inputenc}
\usepackage[british]{babel}

\usepackage[backend=biber,bibstyle=ieee,citestyle=numeric-comp,sorting=none,url=true,doi=true,isbn=false,maxnames=99,minnames=99]{biblatex}
\DefineBibliographyStrings{english}{%
  bibliography = {References},}

\usepackage{amsmath,amsfonts}
\usepackage{url}
\usepackage{xurl}
\usepackage{float}
\usepackage{graphicx}
\graphicspath{{figs/}}
\usepackage{array}
\usepackage{multirow,array}
\usepackage{subcaption}
\usepackage{hyperref}
\hypersetup{colorlinks=true, linkcolor=black, citecolor=black, urlcolor=black, filecolor=black}
\usepackage{microtype}
\usepackage{fancyhdr}
\usepackage{csquotes}
\usepackage{color}
\usepackage{ragged2e}
\usepackage{titlesec}
\usepackage{booktabs}
\usepackage{longtable}
\usepackage{tabularx}
\newcolumntype{L}[1]{>{\RaggedRight\arraybackslash}p{#1}}
\newcolumntype{Y}{>{\RaggedRight\arraybackslash}X}
\usepackage{enumitem}
\usepackage{placeins}
\usepackage[font=small, labelfont=normalfont, skip=4pt]{caption}

\titleformat{\section}[hang]{\normalfont\normalsize\bfseries\RaggedRight}{\thesection}{0.8em}{}
\titleformat{\subsection}[hang]{\normalfont\normalsize\bfseries\RaggedRight}{\thesubsection}{0.8em}{}
\titlespacing*{\section}{0pt}{1.6ex plus 0.4ex}{0.7ex plus 0.2ex}
\titlespacing*{\subsection}{0pt}{1.2ex plus 0.3ex}{0.5ex plus 0.1ex}

\setlength{\textfloatsep}{8pt plus 2pt minus 4pt}
\setlength{\floatsep}{6pt plus 2pt minus 2pt}
\setlength{\intextsep}{6pt plus 2pt minus 2pt}

\newcommand\pic[1]{(Fig. \ref{#1})}
\newcommand\tab[1]{(Tab. \ref{#1})}

\captionsetup[table]{name={Table}, labelsep=period, position=top, justification=raggedright, singlelinecheck=true}
\DeclareCaptionLabelSeparator{figSep}{.\quad}
\captionsetup[figure]{name={Fig.}, labelsep=period, position=bottom, justification=raggedright, singlelinecheck=false}

\pagestyle{fancy}
\fancyhf{}
\renewcommand{\headrulewidth}{0.3pt}
\renewcommand{\sectionmark}[1]{\markright{\thesection\ #1}}
\fancyhead[L]{\small\nouppercase{\rightmark}}
\fancyhead[R]{\small\thepage}
\fancypagestyle{plain}{\fancyhf{}\renewcommand{\headrulewidth}{0pt}}

% Title-page fonts. A course preamble may renew these; the layout stays the same.
\newcommand{\titlepagetitlefont}{\LARGE\bfseries}
\newcommand{\titlepagesubtitlefont}{\large}
\newcommand{\titlepagemetafont}{\normalsize}
\newcommand{\titlepageteamfont}{\normalsize\bfseries}
\newcommand{\titlepagetablefont}{\normalsize}
\newcommand{\titlepagekeyfont}{\small\itshape}

% One row of the who-did-what table on the title page.
\newcommand{\member}[2]{#1 & #2 \\ \addlinespace[3pt]}

% One-page title: title, subtitle, course/date, team, members+contributions, citation key.
% Arguments: title, subtitle, team name, \member{Name}{what they did} rows, course line, date.
% Environment body: citation-key paragraph (wording from the course Agent.md).
\newenvironment{coursereporttitle}[6]{%
  \pagenumbering{Alph}%
  \begin{titlepage}
    \centering
    \vspace*{1.5cm}
    {\titlepagetitlefont #1\par}
    \vspace{0.55em}
    {\titlepagesubtitlefont #2\par}
    \vspace{1.1em}
    {\titlepagemetafont #5\par}
    \vspace{0.25em}
    {\titlepagemetafont Submitted #6\par}
    \vfill
    {\titlepageteamfont Team: #3\par}
    \vspace{0.55em}
    {\titlepagetablefont
    \begin{tabular}{@{}>{\RaggedRight\arraybackslash}p{0.30\textwidth}>{\RaggedRight\arraybackslash}p{0.58\textwidth}@{}}
      \toprule
      \textbf{Member} & \textbf{Contribution} \\
      \midrule
      #4
      \bottomrule
    \end{tabular}\par}
    \vfill
    \begin{minipage}{0.88\textwidth}
      \centering\titlepagekeyfont\ignorespaces
}{%
      \par
    \end{minipage}
    \vspace{1.4cm}
  \end{titlepage}%
  \pagenumbering{arabic}%
}


\usepackage{tempora} % Times New Roman alike font

\usepackage{geometry}
\geometry{left=2.2cm, right=2.2cm, top=1.8cm, bottom=1.8cm, headheight=15pt, headsep=6mm, footskip=10mm}
\usepackage{setspace}
\setstretch{1.15}
\usepackage{indentfirst}
\parindent=1em
\setlength{\parskip}{0pt}

\usepackage[T1]{fontenc}
\usepackage[utf8]{inputenc}
\usepackage[british]{babel}

\usepackage[backend=biber,bibstyle=ieee,citestyle=numeric-comp,sorting=none,url=true,doi=true,isbn=false,maxnames=99,minnames=99]{biblatex}
\DefineBibliographyStrings{english}{%
  bibliography = {References},}

\usepackage{amsmath,amsfonts}
\usepackage{url}
\usepackage{xurl}
\usepackage{float}
\usepackage{graphicx}
\graphicspath{{figs/}}
\usepackage{array}
\usepackage{multirow,array}
\usepackage{subcaption}
\usepackage{hyperref}
\hypersetup{colorlinks=true, linkcolor=black, citecolor=black, urlcolor=black, filecolor=black}
\usepackage{microtype}
\usepackage{fancyhdr}
\usepackage{csquotes}
\usepackage{color}
\usepackage{ragged2e}
\usepackage{titlesec}
\usepackage{booktabs}
\usepackage{longtable}
\usepackage{tabularx}
\newcolumntype{L}[1]{>{\RaggedRight\arraybackslash}p{#1}}
\newcolumntype{Y}{>{\RaggedRight\arraybackslash}X}
\usepackage{enumitem}
\usepackage{placeins}
\usepackage[font=small, labelfont=normalfont, skip=4pt]{caption}

\titleformat{\section}[hang]{\normalfont\normalsize\bfseries\RaggedRight}{\thesection}{0.8em}{}
\titleformat{\subsection}[hang]{\normalfont\normalsize\bfseries\RaggedRight}{\thesubsection}{0.8em}{}
\titlespacing*{\section}{0pt}{1.6ex plus 0.4ex}{0.7ex plus 0.2ex}
\titlespacing*{\subsection}{0pt}{1.2ex plus 0.3ex}{0.5ex plus 0.1ex}

\setlength{\textfloatsep}{8pt plus 2pt minus 4pt}
\setlength{\floatsep}{6pt plus 2pt minus 2pt}
\setlength{\intextsep}{6pt plus 2pt minus 2pt}

\newcommand\pic[1]{(Fig. \ref{#1})}
\newcommand\tab[1]{(Tab. \ref{#1})}

\captionsetup[table]{name={Table}, labelsep=period, position=top, justification=raggedright, singlelinecheck=true}
\DeclareCaptionLabelSeparator{figSep}{.\quad}
\captionsetup[figure]{name={Fig.}, labelsep=period, position=bottom, justification=raggedright, singlelinecheck=false}

\pagestyle{fancy}
\fancyhf{}
\renewcommand{\headrulewidth}{0.3pt}
\renewcommand{\sectionmark}[1]{\markright{\thesection\ #1}}
\fancyhead[L]{\small\nouppercase{\rightmark}}
\fancyhead[R]{\small\thepage}
\fancypagestyle{plain}{\fancyhf{}\renewcommand{\headrulewidth}{0pt}}

% Title-page fonts. A course preamble may renew these; the layout stays the same.
\newcommand{\titlepagetitlefont}{\LARGE\bfseries}
\newcommand{\titlepagesubtitlefont}{\large}
\newcommand{\titlepagemetafont}{\normalsize}
\newcommand{\titlepageteamfont}{\normalsize\bfseries}
\newcommand{\titlepagetablefont}{\normalsize}
\newcommand{\titlepagekeyfont}{\small\itshape}

% One row of the who-did-what table on the title page.
\newcommand{\member}[2]{#1 & #2 \\ \addlinespace[3pt]}

% One-page title: title, subtitle, course/date, team, members+contributions, citation key.
% Arguments: title, subtitle, team name, \member{Name}{what they did} rows, course line, date.
% Environment body: citation-key paragraph (wording from the course Agent.md).
\newenvironment{coursereporttitle}[6]{%
  \pagenumbering{Alph}%
  \begin{titlepage}
    \centering
    \vspace*{1.5cm}
    {\titlepagetitlefont #1\par}
    \vspace{0.55em}
    {\titlepagesubtitlefont #2\par}
    \vspace{1.1em}
    {\titlepagemetafont #5\par}
    \vspace{0.25em}
    {\titlepagemetafont Submitted #6\par}
    \vfill
    {\titlepageteamfont Team: #3\par}
    \vspace{0.55em}
    {\titlepagetablefont
    \begin{tabular}{@{}>{\RaggedRight\arraybackslash}p{0.30\textwidth}>{\RaggedRight\arraybackslash}p{0.58\textwidth}@{}}
      \toprule
      \textbf{Member} & \textbf{Contribution} \\
      \midrule
      #4
      \bottomrule
    \end{tabular}\par}
    \vfill
    \begin{minipage}{0.88\textwidth}
      \centering\titlepagekeyfont\ignorespaces
}{%
      \par
    \end{minipage}
    \vspace{1.4cm}
  \end{titlepage}%
  \pagenumbering{arabic}%
}


\usepackage{tempora} % Times New Roman alike font

\usepackage{geometry}
\geometry{left=2.2cm, right=2.2cm, top=1.8cm, bottom=1.8cm, headheight=15pt, headsep=6mm, footskip=10mm}
\usepackage{setspace}
\setstretch{1.15}
\usepackage{indentfirst}
\parindent=1em
\setlength{\parskip}{0pt}

\usepackage[T1]{fontenc}
\usepackage[utf8]{inputenc}
\usepackage[british]{babel}

\usepackage[backend=biber,bibstyle=ieee,citestyle=numeric-comp,sorting=none,url=true,doi=true,isbn=false,maxnames=99,minnames=99]{biblatex}
\DefineBibliographyStrings{english}{%
  bibliography = {References},}

\usepackage{amsmath,amsfonts}
\usepackage{url}
\usepackage{xurl}
\usepackage{float}
\usepackage{graphicx}
\graphicspath{{figs/}}
\usepackage{array}
\usepackage{multirow,array}
\usepackage{subcaption}
\usepackage{hyperref}
\hypersetup{colorlinks=true, linkcolor=black, citecolor=black, urlcolor=black, filecolor=black}
\usepackage{microtype}
\usepackage{fancyhdr}
\usepackage{csquotes}
\usepackage{color}
\usepackage{ragged2e}
\usepackage{titlesec}
\usepackage{booktabs}
\usepackage{longtable}
\usepackage{tabularx}
\newcolumntype{L}[1]{>{\RaggedRight\arraybackslash}p{#1}}
\newcolumntype{Y}{>{\RaggedRight\arraybackslash}X}
\usepackage{enumitem}
\usepackage{placeins}
\usepackage[font=small, labelfont=normalfont, skip=4pt]{caption}

\titleformat{\section}[hang]{\normalfont\normalsize\bfseries\RaggedRight}{\thesection}{0.8em}{}
\titleformat{\subsection}[hang]{\normalfont\normalsize\bfseries\RaggedRight}{\thesubsection}{0.8em}{}
\titlespacing*{\section}{0pt}{1.6ex plus 0.4ex}{0.7ex plus 0.2ex}
\titlespacing*{\subsection}{0pt}{1.2ex plus 0.3ex}{0.5ex plus 0.1ex}

\setlength{\textfloatsep}{8pt plus 2pt minus 4pt}
\setlength{\floatsep}{6pt plus 2pt minus 2pt}
\setlength{\intextsep}{6pt plus 2pt minus 2pt}

\newcommand\pic[1]{(Fig. \ref{#1})}
\newcommand\tab[1]{(Tab. \ref{#1})}

\captionsetup[table]{name={Table}, labelsep=period, position=top, justification=raggedright, singlelinecheck=true}
\DeclareCaptionLabelSeparator{figSep}{.\quad}
\captionsetup[figure]{name={Fig.}, labelsep=period, position=bottom, justification=raggedright, singlelinecheck=false}

\pagestyle{fancy}
\fancyhf{}
\renewcommand{\headrulewidth}{0.3pt}
\renewcommand{\sectionmark}[1]{\markright{\thesection\ #1}}
\fancyhead[L]{\small\nouppercase{\rightmark}}
\fancyhead[R]{\small\thepage}
\fancypagestyle{plain}{\fancyhf{}\renewcommand{\headrulewidth}{0pt}}

% Title-page fonts. A course preamble may renew these; the layout stays the same.
\newcommand{\titlepagetitlefont}{\LARGE\bfseries}
\newcommand{\titlepagesubtitlefont}{\large}
\newcommand{\titlepagemetafont}{\normalsize}
\newcommand{\titlepageteamfont}{\normalsize\bfseries}
\newcommand{\titlepagetablefont}{\normalsize}
\newcommand{\titlepagekeyfont}{\small\itshape}

% One row of the who-did-what table on the title page.
\newcommand{\member}[2]{#1 & #2 \\ \addlinespace[3pt]}

% One-page title: title, subtitle, course/date, team, members+contributions, citation key.
% Arguments: title, subtitle, team name, \member{Name}{what they did} rows, course line, date.
% Environment body: citation-key paragraph (wording from the course Agent.md).
\newenvironment{coursereporttitle}[6]{%
  \pagenumbering{Alph}%
  \begin{titlepage}
    \centering
    \vspace*{1.5cm}
    {\titlepagetitlefont #1\par}
    \vspace{0.55em}
    {\titlepagesubtitlefont #2\par}
    \vspace{1.1em}
    {\titlepagemetafont #5\par}
    \vspace{0.25em}
    {\titlepagemetafont Submitted #6\par}
    \vfill
    {\titlepageteamfont Team: #3\par}
    \vspace{0.55em}
    {\titlepagetablefont
    \begin{tabular}{@{}>{\RaggedRight\arraybackslash}p{0.30\textwidth}>{\RaggedRight\arraybackslash}p{0.58\textwidth}@{}}
      \toprule
      \textbf{Member} & \textbf{Contribution} \\
      \midrule
      #4
      \bottomrule
    \end{tabular}\par}
    \vfill
    \begin{minipage}{0.88\textwidth}
      \centering\titlepagekeyfont\ignorespaces
}{%
      \par
    \end{minipage}
    \vspace{1.4cm}
  \end{titlepage}%
  \pagenumbering{arabic}%
}

% MSD case-report format. Taken from Satera-C1-Report.docx; CASE-RULES
% require 11-12 pt and a 5-page body (title page, references, appendix excluded).
% Every MSD assignment: 11 pt letter, common preamble, then this file.

\geometry{letterpaper, left=0.945in, right=0.945in, top=0.866in, bottom=0.866in, headheight=13pt, headsep=4.2mm, footskip=8mm}

% Calibri-metric font (docx body is Calibri 11 pt).
\usepackage[sfdefault]{carlito}
\renewcommand{\familydefault}{\sfdefault}

% Normal: 11 pt, line 259/240, 5 pt after each paragraph, no first-line indent.
\setstretch{1.079}
\setlength{\parindent}{0pt}
\setlength{\parskip}{5pt}
\raggedbottom

% Heading 1: 13 pt bold, 10 pt before, 4 pt after.
\titleformat{\section}[hang]{\fontsize{13}{16}\selectfont\bfseries\raggedright}{\thesection}{0.7em}{}
\titlespacing*{\section}{0pt}{10pt}{4pt}

% Heading 2: 11.5 pt bold, 6 pt before, 3 pt after.
\titleformat{\subsection}[hang]{\fontsize{11.5}{14}\selectfont\bfseries\raggedright}{\thesubsection}{0.7em}{}
\titlespacing*{\subsection}{0pt}{6pt}{3pt}

% Caption style in the docx is 9 pt.
\captionsetup{font={footnotesize}, labelfont={bf}, skip=4pt}

% Word-style header: page number only.
\fancyhf{}
\renewcommand{\headrulewidth}{0pt}
\fancyhead[R]{\fontsize{11}{13}\selectfont\thepage}

% Hanging references (docx hanging indent 454 twips ~ 0.8 cm).
\setlength{\bibhang}{0.8cm}
\renewcommand*{\bibfont}{\fontsize{11}{13}\selectfont}

% Title page: same layout as common (team, members, who did what).
% Sizes from Satera-C1-Report.docx.
\renewcommand{\titlepagetitlefont}{\fontsize{22}{26}\selectfont\bfseries}
\renewcommand{\titlepagesubtitlefont}{\fontsize{13}{16}\selectfont}
\renewcommand{\titlepagemetafont}{\fontsize{11}{14}\selectfont}
\renewcommand{\titlepageteamfont}{\fontsize{12}{15}\selectfont\bfseries}
\renewcommand{\titlepagetablefont}{\fontsize{11}{14}\selectfont}
\renewcommand{\titlepagekeyfont}{\fontsize{9.5}{12}\selectfont\itshape}

\bibliography{ref.bib}

\begin{document}

\begin{coursereporttitle}
  {The Satera Team at Imatron Systems}
  {Case C1 report --- getting Satera back on track in 30 days}
  {Temporary name}
  {% Who did what (CASE-RULES, title slide/page). Source: team chat, 8--18 Sep 2026.
    \member{Lishchenko Ivan}{Formed and registered the team; organised the pre-case call; collected the L1a concepts to reference; drafted the assumption on combining both engineers' strengths.}
    \member{Smirnov Fedor}{Proposed the first recommendation and reasons; built the first slide deck; proposed the leadership-style root problem; rewrote the risks and added the fourth reason (continuous intervention); finalised the slides.}
    \member{Cherkashin Nikita}{Wrote the draft-deck reason on the engineers' complementary strengths (Lovas's detail focus, Bennett's optimisation) and why a collaboration gives both leaders a significant role.}
    \member{Faki Doosuur Doris}{Drafted reasons 3 and 4 for the case-day draft; relayed the lecturer's critique of our root problem and restarted the analysis from it; proposed the affiliative-versus-authoritative root and coaching for Bennett; edited the revised slides.}
    \member{Gadzama Ezekiel John}{Drafted the revised recommendation (Pinto's own call); merged the two revised slide versions into one deck; added case page references; compiled the report and set it up in LaTeX.}
  }
  {Managing Software Development \textperiodcentered\ Innopolis University \textperiodcentered\ Fall 2026}
  {18 September 2026}
Citation key. ``p.N'' is a page of the case \parencite{amabile2003satera} and ``Exhibit~N'' a case exhibit. ``L1a s.N'' and ``L1b s.N'' are slide numbers in the course lectures \parencite{sadovykh2026l1a,sadovykh2026l1b}. [A] marks an assumption: the case is silent, so it cites nothing. Appendix~A lists what each assumption rests on and what changes if it is wrong.
\end{coursereporttitle}

\section{Recommendation}
\textbf{Recommendation.}
Gary Pinto should switch from the affiliative style he defaults to and take the support-structure decision with authority: announce the honeycomb as his own call this week, ask Bennett to make it lighter, and hold short, regular check-ins so that disagreements between the two senior engineers are dealt with early.

\textbf{What changes?}
The mechanical design, which cannot be finalised until the support structure is chosen (p.~1), restarts. Because the choice is Pinto's rather than a team vote, neither engineer is seen to lose to ``the team''. Bennett takes on the weight reduction the honeycomb needs [A], and his wheel is documented as an option for future work, which suits a project meant to be a platform technology for future products (p.~1). Fast intervention keeps the arguments from escalating again [A].

\textbf{Two outputs.}
L1b s.36 asks that a decision name two outputs: the solution and the risk accepted with it. Our solution is the honeycomb, decided by Pinto. The risk we assume is that Bennett does not accept a supporting role, or that the honeycomb stays heavier than Pinto would like (p.~8). Section~5 deals with both.

\section[Root problem]{Root problem: Pinto leads with an affiliative style where the decision needs authority}
L1a s.28 describes six leadership styles, each with a situation it fits and a cost. The affiliative style (``people come first'') fits repairing trust or high stress; its cost is that hard feedback goes unsaid. The case shows this is Pinto's default. He builds team spirit by reading team members' bios aloud and pairs praise with constructive criticism (p.~5). He handles conflict ``generally through discussion'' rather than discipline, which worked on a previous team but has failed with Lovas and Bennett (p.~6). He made the design choice depend on the whole team's reaction (p.~6) and closed the demonstration by asking everyone to think it over (p.~8), although he had already made his choice (p.~2) and did not know how to tell them (p.~3). The hard feedback --- that Bennett's design will not be pursued --- is exactly what has gone unsaid: Bennett still hopes Pinto will give his design a chance (p.~9).

The situation needs the opposite. The prototype is due on 1 April 2000 under a contract that cuts payments for delay, the mechanical design phase is already late (p.~1), and Levinger has asked Pinto for quick action (p.~1--2). For a decision that sets direction, L1a s.29 prescribes democratic input followed by an authoritative ending. Pinto has the input --- two prototypes, a vibration test and the team's e-mails (p.~6--8) --- but has not supplied the ending. The lecture says the authoritative style fails when the leader lacks credibility among experts (L1a s.28). That condition does not apply: Lovas relies on Pinto for complicated tests, Bennett is impressed by his grasp of the design (p.~6), and management regards him as an effective leader (p.~5). The style he needs is one he can use.

\textbf{Why this is the root and not a symptom.}
Every other problem the case describes follows from this one. The support structure is undecided, so the mechanical design is stuck (p.~1). While the choice is open, each engineer keeps arguing for his own design: the day after the test they argue again, and Bennett still hopes to win (p.~9). Baxter writes that the team freezes whenever a major design decision is due, and Pinto knows the person avoiding it is him (p.~2). Team members were too intimidated to state a view in the meeting (p.~3), and nobody stepped into the debate (p.~7). The reverse test points the same way. Attempts to fix the relationship without changing how decisions are made have already failed: team building at the start (p.~4), Pinto's usual conflict discussions (p.~6), and two weeks of working apart, which produced two competing prototypes rather than a decision (p.~2).

We do not treat the conflict itself as the root. The case says both engineers are normally very professional and differ in how they approach the work (p.~2); L1a s.12--13 finds that personality does not predict team performance, and in the largest study L1b cites, relationship conflict did not reduce team effectiveness (L1b s.16). L1a s.34 also warns that the six styles rest on unpublished, correlational work with executives, so we do not let that model carry the argument alone: the reasons in Section~3 rest equally on the decision and team findings in L1b.

\section[Why the recommendation follows]{Why the recommendation follows}
\subsection{Nothing moves until someone closes the decision}
The mechanical design cannot be finalised until the support structure is complete (p.~1). Both designs are feasible, but one is much more likely to perform at launch, and the team saw this too (p.~3). Pinto cannot remove either engineer: he needs both, and their relationships with the customer, through prototype development (p.~2). So the decision must be one both can keep working under. L1b s.23 presents analysis paralysis as: not deciding badly, but not deciding at all, because every option has drawbacks. Satera's options fit this pattern exactly. The wheel is light but began to warp within 20 seconds (p.~7); the honeycomb held for three minutes but may be too heavy (p.~8). Waiting for an option without drawbacks is waiting indefinitely.

We assume Satera had no agreed criteria for accepting ideas [A], since the case describes none; L1b s.29 says a group should agree how it will decide before it decides, which is what the week-1 rules in Section~4 do. We also assume Bennett's strength is the missing piece, not a threat [A]: the honeycomb's weakness is weight (p.~8), Bennett's structure was impressively light (p.~7), and Pinto wants the two to make the chosen structure lighter together (p.~3).

\subsection{A leader's call protects Bennett; a team call isolates him}
Satera's members report considerably less work-group support than other RID teams (p.~4; Exhibit~3, a single seven-point item on a daily survey). Rowling says the team never built the trust needed to tolerate its diversity, and that several members ignored the team-building effort made at the start (p.~4). In a team like this, a decision that looks like the team's verdict splits it into camps. The team has already sided with the honeycomb, privately, by e-mail (p.~8), and Bennett already doubts that Satera or ISI will ever trust his ingenuity (p.~8).

L1b s.28 allows two ways for a group to decide. Majority rule on a complex, high-stakes decision often produces a win/lose result --- here, Bennett against the colleagues he must work with until the prototype ships. Consensus is not available: the two engineers have failed to agree for weeks (p.~2), and the deadline fails the lecture's timeliness criterion. That leaves the leader's ending that L1a s.29 prescribes. Announced as Pinto's call, with his reasons, the decision is one Bennett lost to his manager's judgement rather than to his colleagues. We assume that affiliative gestures without direction do not build real trust [A], which matches L1b s.9: trust grows through reliability and transparency, not team-building events.

\subsection{Continuous intervention lets the situation stabilise}
The debate between the two engineers has turned into unproductive bickering on an increasingly regular basis (p.~1), and after the demonstration it turned heated again within a day (p.~9). Team building happened once, at the start (p.~4), and Pinto's usual discussion method has already failed with this pair (p.~6). One announcement will settle the structure but not every later disagreement about how to lighten it. L1b s.11 advises putting a contested decision in writing with an owner and a date, and treating the first real argument as expected rather than something to smooth over. We assume that short, regular check-ins by Pinto keep the bickering from escalating [A], because the case shows it grows when left alone (p.~1, 9).

\section{The 30-day plan}
Pinto owns every step; the tiger team checks the measures and reports to Levinger at day 30. The case gives no weight figure, so the target is set in week 1 [A].

\begin{table}[H]
\centering
\small
\caption{30-day plan}
\label{tab:plan}
\begin{tabularx}{\textwidth}{l L{2.3cm} X X}
\toprule
\textbf{When} & \textbf{Who} & \textbf{Action} & \textbf{How we know it worked} \\
\midrule
Day 1 & Pinto & Short meeting: states the honeycomb decision as his own call and gives his reasons (test results, p.~7--8). Asks Bennett to lead the weight reduction; Lovas continues the structure. & Both engineers leave with a named task; nobody reopens the choice. \\
Week 1 & Pinto, with Lovas and Bennett & Sets three written rules: idea time (a fixed slot where new ideas are heard), stated reasons (every objection comes with one), escalation (a disagreement not settled in 48 hours goes to Pinto, who decides). & Rules written and shared by day 5; weight target agreed. \\
End of week 1 & Pinto & Full-team meeting: restates the decision and reasons, credits both designs, records the wheel as a future platform option (p.~1). & Mechanical design work has resumed. \\
Weeks 2--4 & Pinto & Short check-ins with Lovas and Bennett, twice a week, to catch new conflict early. & No design question open past the 48-hour escalation window. \\
Week 2 & Bennett, Lovas & First weight-reduction result, retested on the vibration table. & Lighter than the original honeycomb and still passing the test (p.~8); if not, risk 2 applies. \\
Day 30 & Pinto, tiger team & Review with Levinger. & Design resumed, weight target met, no new blockers. \\
\bottomrule
\end{tabularx}
\end{table}

\section{What could go wrong}
\textbf{Risk 1 (biggest): Bennett will not accept a supporting role.}
This is the assumption the plan depends on most [A]. Against it: Bennett has said the team would progress faster if he were sole leader of the mechanical work (p.~5), and he appears to lose a piece of himself each time an idea is shelved (p.~9). For it: he respects Pinto's grasp of the design (p.~6), he acknowledged his structure's endurance problem during the test (p.~7), and the team is intrinsically motivated (p.~3). Expectancy theory (L1a s.24) explains why the weight task matters: motivation needs effort to lead to results and results to a reward worth having, and one belief at zero makes the whole chain zero. Owning a problem that visibly reaches the product restores that link; responsibility is also one of Herzberg's motivators (L1a s.22). If the assumption is wrong, Bennett disengages and Satera loses half its senior mechanical leadership, which Pinto says he cannot afford (p.~2). The fallback is one-to-one coaching with Bennett (L1a s.28), which works only if he is willing --- the lecture's own condition for coaching.

\textbf{Risk 2: the honeycomb shows no weight reduction by week 2.}
[A] The fallback from our presentation is to switch to Bennett's prototype. We keep it with a condition, because the wheel warped within 20 seconds (p.~7) and Lovas doubted it could be improved without new structural problems (p.~7). We would switch only if Bennett's structure passes the same three-minute vibration test the honeycomb passed (p.~8); otherwise the team accepts the extra weight, as its members already did when they chose the honeycomb (p.~8).

\textbf{Risk 3: Pinto cannot act decisively for a prolonged period.}
[A] The choice itself was not difficult, yet Pinto has still not told the team, which is ``holding its breath'' (p.~2--3). The fallback is that the tiger team takes charge of delivering decisions while Pinto stays as adviser and trusted figure. We keep Pinto in the team because L1b s.12 finds that re-composing a team costs respect and causes anger in discussions, and both engineers rely on him (p.~6).

\section[What changed]{What changed since the presentation}
In class we named the root problem as a personal conflict between the two senior engineers and recommended that they work on each other's prototypes for two to three weeks. In cross-examination the lecturer questioned how we chose the root. We changed it for three reasons. The case describes a disagreement about whose design is used, between two normally professional engineers (p.~2), not a personal feud. Fixes aimed at the relationship had already been tried and failed (p.~2, 4, 6). And swapping prototypes spends two to three weeks of our 30 days and still ends without a decision, which L1b s.23 names as the common failure, under a contract that penalises delay (p.~1). The leadership-style diagnosis explains both the conflict and the deadlock, and the recommendation now follows from it. We also added the condition on the risk-2 fallback after rereading the test results.

\nocite{msd2026c1,goleman2000,herzberg1968,soomro2016,thompson2015,hoffmann2021,verwijs2024,hoover2010,rebori1998}
\clearpage
\printbibliography

\clearpage
\section*{Appendix A --- Assumptions}
Each assumption rests on the case's silence. The last column says what changes if it is wrong.

\begin{table}[H]
\centering
\small
\caption{Assumptions}
\label{tab:assumptions}
\begin{tabularx}{\textwidth}{Y Y Y}
\toprule
\textbf{Assumption [A]} & \textbf{What it rests on} & \textbf{If it is wrong} \\
\midrule
Bennett accepts a supporting role (load-bearing) & He respects Pinto (p.~6) and admitted his design's flaw (p.~7) & He disengages; coaching fallback, then the risk-1 plan (Section~5) \\
Satera had no agreed criteria for accepting ideas & The case describes none & Then Pinto did not apply them, which strengthens the case for his deciding \\
Bennett's strength is the missing piece & Honeycomb's weakness is weight (p.~8); wheel was light (p.~7) & Lovas lightens it alone; Bennett needs another task he owns \\
A leader's call isolates Bennett less than a team vote & Team sided privately (p.~8); low support (p.~4) & Bennett resents Pinto instead; handled as risk 1 \\
Affiliative gestures without direction do not build trust & Kickoff team building ignored (p.~4); L1b s.9 & More team building would help, and support ratings would rise \\
Regular check-ins keep the bickering from escalating & Bickering grows when left alone (p.~1, 9) & Check-ins become another arena; Pinto decides disputes in writing instead \\
Fast intervention keeps arguments from escalating & Same evidence as above & As above \\
A weight target can be set in week 1 & The case gives no figure & Measure against the original honeycomb \\
Pinto can act decisively once he chooses to & Management rates him effective (p.~5) & Risk-3 fallback \\
\bottomrule
\end{tabularx}
\end{table}

\end{document}
